\chapter{Tập con Common Command Code}
\label{app:ccc}

Phụ lục này chi tiết hóa tập con CCC được thiết kế hỗ trợ, bổ trợ cho \cref{sec:bg-ccc}. Các mã được lấy
trực tiếp từ kiểu liệt kê \texttt{i3c\_ccc\_e} trong \texttt{i3c\_pkg}. Địa chỉ broadcast dành riêng của
I3C là \texttt{7'h7E}, dùng để mở đầu các CCC Broadcast và ENTDAA.

\begin{longtable}{L{2.8cm} C{2.0cm} C{2.2cm} L{6.0cm}}
  \toprule
  \textbf{CCC} & \textbf{Mã} & \textbf{Loại} & \textbf{Mô tả} \\
  \midrule\endhead
  \texttt{ENEC}      & 0x00 & Broadcast & Enable Events Command — bật báo sự kiện cho tất cả Target. \\
  \texttt{DISEC}     & 0x01 & Broadcast & Disable Events Command — tắt báo sự kiện cho tất cả Target. \\
  \texttt{ENTDAA}    & 0x07 & Broadcast & Enter Dynamic Address Assignment — khởi động quy trình gán địa
                                          chỉ động. \\
  \texttt{DIR\_ENEC}  & 0x80 & Direct    & Enable Events tới một Target theo Dynamic Address. \\
  \texttt{DIR\_DISEC} & 0x81 & Direct    & Disable Events tới một Target theo Dynamic Address. \\
  \bottomrule
\end{longtable}

Trong hiện thực, một CCC Direct được nhận biết qua bit cao nhất của mã lệnh (\texttt{cmd[7] = 1}); máy
trạng thái \texttt{flow\_active} dùng bit này ở trạng thái \texttt{IssueImmediateCcc} để chọn đúng trình tự
pha và điều kiện dừng giữa Direct và Broadcast (xem \cref{sec:flow-merge}).
